\subsection{Bounding the average cusp correction by counting}
\label{sec:cusp-counting-mean}

The averaging estimate has a short proof once the original quotient
measure and counting identity are identified. Write $\log^+x=\max(0,\log x)$.

\begin{lemma}\label{lem:cusp-counting-tail}
On a probability space let $\lambda>0$ and $N_t\ge0$ be measurable,
and fix $r>0$ and $0<a<k$. Suppose $N_r$ is integrable and, for
$0<t\le r$,
\[
 \mathbf1_{\{\lambda<t\}}\le N_t,\qquad
 N_r\le(3r/\lambda)^a\mathbf1_{\{\lambda<r\}},\qquad
 \int N_t\le(t/r)^k\int N_r.
\]
Then $\log^+(r/\lambda)$ is integrable and
\[
 \int\log^+(r/\lambda)
 \le\frac{a\log3+\log(k/(k-a))}{k-a}+\frac1k.
\]
For $a\ge2$ and $k\ge3a$, this is at most $23/24<1$.
\end{lemma}

\begin{proof}
Put $M=\int N_r$. The hypotheses give
$\Pr(\lambda<t)\le\min(1,M(t/r)^k)$. If $M\ge1$, integrating this
tail, with the split at $u=M^{a/k}$, gives
\[
 M\le3^a\int(r/\lambda)^a\mathbf1_{\{\lambda<r\}}
 \le3^a\left(1+\int_1^\infty\min(1,Mu^{-k/a})\,du\right)
 =\frac{3^ak}{k-a}M^{a/k}.
\]
Thus $M\le M_*=(3^ak/(k-a))^{k/(k-a)}$, also when $M<1$.
Tonelli's theorem now gives
\[
 \int\log^+(r/\lambda)
 \le\int_0^\infty\min(1,M_*e^{-ks})\,ds
 =\frac{\log M_*+1}{k}.
\]
This proves integrability as well as the claimed bound. For $a\ge2$,
$k\ge3a$, it is at most
$\frac12\log3+\frac14\log(3/2)+\frac16\le23/24$;
use $\log3\le4/3$ and $\log(3/2)\le1/2$.
\end{proof}

Here is how the pointwise hypotheses arise. Let $E$ be CM of degree
$d=2f$, let $L$ be an $\mathcal O_E$-stable full lattice in a Hermitian
$E$-space of Witt index at most $j$, and let $\lambda$ be its shortest
nonzero vector length in the original Hodge norm. Suppose
$b(v,w)=\operatorname{Tr}_{E/\mathbb Q}(\xi h(v,w))$ is integral on
$L$, with $\xi\ne0$, and
\[
 |b(v,w)|\le\|v\|\|w\|,\qquad
 \|\alpha v\|\le\max_\sigma|\sigma\alpha|\,\|v\|.
\]
Choose integral elements spanning $E$ over $\mathbb Q$, all of whose
conjugates have absolute value at most $v_r\ge1$, and put $r=v_r^{-1/2}$.
For $\|v\|,\|w\|<r$, the integers $b(\alpha v,w)$ have absolute
value less than one. They vanish; the spanning property and the trace
pairing imply $h(v,w)=0$. All short vectors therefore span a totally
isotropic $E$-space, of real dimension at most $2fj$. Balls of radius
$\lambda/2$ around distinct vectors are disjoint. Comparing their
volumes inside the ball of radius $t+\lambda/2$ gives, for the number
$N_t$ of nonzero isotropic vectors of norm less than $t$,
\[
 1\le N_t\le(3t/\lambda)^{2fj}\quad(\lambda<t\le r),
 \qquad N_t=0\quad(t\le\lambda).
\]
Thus the lemma applies with $a=2fj$ and $k=2f(N-1)$ when $N\ge3j+1$,
provided the original quotient probability measure makes $N_r$
integrable and satisfies $\int N_t=(t/r)^k\int N_r$.

For abelian $E/\mathbb Q$, the spanning elements admit a simple bound.
Let $c$ be its conductor. The conductor--discriminant formula gives
$c\le\operatorname{rd}(E)^2$: at each prime dividing $c$, the
characters with smaller conductor exponent form a proper subgroup,
so at least half the $d$ characters have the maximal exponent.
The integral elements
\[
 \operatorname{Tr}_{\mathbb Q(\zeta_c)/E}(\zeta_c^i),
 \qquad0\le i<\varphi(c),
\]
span $E$ over $\mathbb Q$ and have all conjugates bounded by
$v_r=\varphi(c)/d$. This is a rational spanning assertion; their
integral span may have nontrivial index in $\mathcal O_E$.
The classical conductor identities and cyclotomic inclusion are in
Milne, \emph{Class Field Theory}, v4.03,
\href{https://www.jmilne.org/math/CourseNotes/CFT.pdf}{Theorems I.4.16
and V.3.27, and Example V.3.10}.

Assume the same measurable, nonnegative cusp correction $U$ satisfies
the original pointwise comparison
\[
 U\le fj\bigl(2\log^+(1/\lambda)+\log(2fj\operatorname{rd}(E))\bigr).
\]
Since $\log^+(1/\lambda)\le\log^+(r/\lambda)+\tfrac12\log v_r$,
the lemma proves its integrability and
\[
 \int U\le fj\bigl(2+\log v_r+\log(2fj\operatorname{rd}(E))\bigr)
 \le fj\bigl(2+\log j+3\log\operatorname{rd}(E)\bigr).
\]
For $B\ge4fN^2$ and $j\ge1$, one has $B\ge64j$; hence, whenever
$\Lambda_B\ge\log B+\log\operatorname{rd}(E)$,
\[
                         \int U\le3fj\Lambda_B.
\]
This supplies the numerical allowance $C_u=3$. Application to the
arithmetic inequality still requires the original lattice and
covolume, the normalized quotient measure and its counting identity,
and the identification of this same integral with the signed cusp
term. The arithmetic model and the one-sided boundary limit are
retained hypotheses.
